\documentclass[../DoAn.tex]{subfiles}
\begin{document}

\begin{center}
    \Large{\textbf{ABSTRACT}}\\
\end{center}
\vspace{1cm}
Virtual Network Embedding (VNE) is a fundamental problem in network virtualization: each virtual network request must be mapped onto a shared physical substrate so as to satisfy resource constraints while optimizing service efficiency. VNE is NP-Hard, and real deployments are harder still due to their online nature and large scale. Heuristic methods are fast but rely on expert knowledge and generalize poorly; recent reinforcement-learning (RL) methods overcome this but tend to be unstable, sample-inefficient, and ignore readily available topological knowledge.

This thesis proposes \textit{CARL-VNE} (Candidate-Anchored Reinforcement Learning for Virtual Network Embedding), a method that combines graph neural networks with reinforcement learning for multi-domain VNE. It uses a \textit{hierarchical} graph encoder (intra-domain GCN combined with an inter-domain GCN) and an attention-based candidate decoder with an explicit cost bias and a hard feasibility mask. Training proceeds in two stages: an imitation-learning warm-start from the MP-VNE heuristic, followed by actor-critic fine-tuning with a Kullback-Leibler anchor. The core contribution is the finding that \textit{candidate selection} is the effective lever for RL, and that making this selection stochastic during training is necessary to surpass the imitation ceiling.

Experiments on the $50$-node scenario show that the proposed method outperforms the MP-VNE heuristic on $19$ of $21$ VNR-distribution scenarios (reaching $32.7\%$ acceptance versus $25.5\%$ at the central scenario), demonstrating a robust advantage across a wide range of operating conditions while converging stably during training.

\end{document}
